\section{问题提出的背景}

\subsection{背景介绍}

\subsubsection{大语言模型与智能体技术的发展现状}

近年来，大语言模型（Large Language Model, LLM）技术实现了跨越式发展\cite{brown2020gpt3,openai2023gpt4}，以 GPT 系列、通义千问、文心一言等为代表的模型，在自然语言理解、逻辑推理、任务规划等方面的能力得到了显著提升，为人工智能从感知智能向决策智能的跨越奠定了基础\cite{brown2020gpt3}。

在此背景下，智能体（Agent）技术成为大模型落地应用的核心方向 —— 通过赋予大模型工具调用、自主规划、多轮迭代的能力，智能体能够突破大模型的知识边界与能力局限，对接外部系统与工具，完成复杂的现实世界任务，成为当前人工智能领域的研究与产业落地热点\cite{yao2023react,schick2023toolformer,qin2023toollearning}。

目前，国内外已形成了成熟的智能体开发基础设施，包括 LangChain、LlamaIndex 等开源框架\cite{langchain2023,llamaindex2023}，提供了标准化的工具封装、任务规划、记忆管理等核心能力，大幅降低了智能体系统的开发门槛。同时，主流大模型均已原生支持函数调用（Function Calling）能力\cite{openai2023gpt4}，能够精准解析用户需求，生成符合格式要求的工具调用指令，为智能体工具调用能力的落地提供了底层支撑\cite{schick2023toolformer,qin2023toollearning}。

\subsubsection{AI for Science 领域的发展与工具需求}

AI for Science（人工智能赋能科学研究）已成为人工智能与基础科学交叉融合的核心赛道，被广泛认为是加速科学发现、提升科研效率的核心技术手段\cite{jumper2021alphafold,butler2018ml4mat}。当前，科学研究场景中存在大量重复性、流程化的工作，包括数值计算、统计分析、实验数据可视化、文献检索、专业数据处理等，这些工作普遍依赖各类专业软件、开源计算库与分析工具，存在工具种类繁多、接口格式不统一、使用门槛高、操作流程繁琐等问题。

对于科研人员而言，尤其是非计算机专业的科研工作者，往往需要花费大量时间学习各类工具的使用方法、处理不同工具之间的数据兼容问题，无法将核心精力聚焦于科学问题本身。因此，科研领域对能够自动化调用各类工具、完成标准化科研流程的智能系统，存在迫切的现实需求。具备工具调用能力的科学智能体，能够有效降低专业工具的使用门槛，实现科研流程的自动化执行，成为 AI for Science 领域的重要研究方向\cite{shen2023hugginggpt,jumper2021alphafold}。

\subsubsection{当前存在的主要问题}

当前主流的智能体框架与工具调用技术，在通用场景中已得到了广泛应用\cite{langchain2023,llamaindex2023}，但在科学研究场景中仍存在明显的适配性缺陷，无法满足科研工作的专业性、严谨性要求，核心问题集中于三个方面：

第一，缺乏面向科学领域的统一工具抽象规范。当前科学计算领域的工具种类繁多，不同开源库、专业软件的接口格式、参数要求、数据返回格式差异极大，通用智能体框架缺乏标准化的工具封装规范，导致工具集成成本高、适配难度大，无法快速扩展支持各类科研工具。

第二，科学任务理解与拆解能力不足。科学研究任务具有专业性强、逻辑严谨、流程固定的特点，通用智能体的任务规划机制无法精准理解专业科研需求\cite{yao2023react,wang2023plansolve}，难以将模糊的科研问题拆解为严谨、可执行的工具调用序列，易出现任务拆解逻辑错误、工具匹配不当等问题。

第三，工具调用的容错性与结果可靠性不足。科研工作对计算结果的精度与严谨性有极高要求，而当前大模型工具调用普遍存在参数幻觉、调用逻辑错误、异常处理能力弱等问题\cite{schick2023toolformer,qin2023toollearning}，缺乏针对科学场景的结果校验机制，无法保障输出结果的准确性与可靠性，难以满足科研工作的合规要求。

\subsection{本研究的意义和目的}

\subsubsection{研究意义}

本研究针对科学场景的核心需求，构建适配科研任务特性的智能体工具调用优化体系，完善垂直领域智能体工具调用的理论框架，弥补现有通用智能体框架在科学场景下的适配性缺陷，为科研领域智能体系统的设计与开发提供可参考的实现范式，丰富 AI for Science 领域的技术体系\cite{jumper2021alphafold,butler2018ml4mat}。

本研究设计的系统能够有效降低科研人员使用专业计算工具的门槛，实现数值计算、统计分析、数据可视化等高频科研任务的自动化执行，大幅减少科研工作中的重复性劳动，提升科研工作效率。同时，系统采用模块化、可扩展的设计思路，能够快速适配生物信息学、材料科学、环境科学等多个细分科研领域的需求\cite{butler2018ml4mat}，具备较强的落地应用潜力，能够为基层科研工作者提供轻量化、易用型的科研辅助工具。

\subsubsection{研究目的}

本研究旨在设计并实现一套面向科学场景的智能体工具调用优化系统，解决当前智能体在科研场景中工具适配难、任务拆解不准、调用容错性差、结果不可靠等核心问题\cite{yao2023react,schick2023toolformer,qin2023toollearning}。通过标准化的工具封装、专项优化的任务拆解机制、轻量化的调用优化与容错校验方案，提升智能体在科学场景中工具调用的准确率、成功率与执行效率，最终实现科研任务的全流程自动化执行，为科研人员提供稳定、易用的智能辅助工具。
